\chapter{General Theory}

\section{Matrix Lie Groups}

\lecture{1}{Tue 20 Jan 2026 10:59}{Syllabus day}

Tuesday 1/27 we have class in cds548 for a seminar on derived symplectic structure.

\begin{definition}\label{1.1.1}
	A Lie Group $G$ is a group object in the category $\mathcal{M} \mathrm{an} $ of smooth manifolds. In particular, it is a group $G$ with a smooth structure that is compatible with the group structure in the sense that multiplication $\mu : G\times G \to G$ is smooth.
\end{definition}

Consider for example the group $G=\operatorname{SU} (2)$ of Hermitian matrices $A A^\dagger=1$. We have $\operatorname{SU} (2)\cong S^3$ as spaces, so there exists a manifold structure on $\operatorname{SU} (2)$. It then becomes a Lie group.

\begin{example}[Matrix groups]\label{1.1.2}
	Consider $\operatorname{GL} (n,\mathbb{R} )\coloneqq \Aut_{\mathbb{R} \text{-} \mathcal{V} \mathrm{ect} }(\mathbb{R} ^n) $. This is an open subset of $\mathbb{R} ^{n^2} $, and thus has an $n^2$-manifold structure. We can also consider $\operatorname{GL} (n,\mathbb{C} )$ and obtain an open subset of $\mathbb{C} ^{n^2} $. This is a subgroup of $\operatorname{GL} (2n,\mathbb{R} )$ in the obvious way. Note that $\operatorname{GL} (n,\mathbb{C} )$ acts on $\mathbb{C} ^n$ by multiplication, and thus acts on $\mathbb{R} ^{2n} $ via the identification $\vec{x}+i \vec{y} \leftrightarrow (\vec{x} , \vec{y} )$.

	More explicitly, take $\mathcal{A} \in \operatorname{GL} (n,\mathbb{C} )$ and $z\in\mathbb{C} ^n$. Then we can write $\mathcal{A} =A+iB$ and $z=x+iy$. Then
	\[
		\mathcal{A} z=\left( Ax+iAy \right) +(-By+iBx)
	.\]
	This can be rewritten of the form
	\[
		\begin{pmatrix}
		A & -B \\
		B & A
		\end{pmatrix}\begin{pmatrix}
		x \\
		y
		\end{pmatrix}
	.\]
	This is equivalent to the statement that the matrix commutes with the matrix
	\[
		J=\begin{pmatrix}
		0 & -I \\
		I & 0
		\end{pmatrix},\qquad I\text{ the $2n$-identity matrix}
	.\]
	We say it \emph{commutes with the complex structure} of multiplication by $i$.
\end{example}

Continuing to work over $\mathbb{C} ^n$, a function is holomorphic when $\partial _{\overline{z} } f=0$, meaning $\frac{1}{2} (\partial _x+i \partial _y)f=0$. The reason there is a plus for $\overline{z} $ is because this is dual to $\mathrm{d} z=\mathrm{d} x+i\mathrm{d} y$ and $\mathrm{d} \overline{z} =\mathrm{d} x-i\mathrm{d} y$ or something.

Using the notion of holomorphy, we can define complex manifolds.

\begin{definition}\label{1.1.3}
	A complex $n$-manifold is a topological manifold $M$ with charts $M \supseteq U \cong V \subseteq \mathbb{C} ^n$ such that the coordinate transformations are holomorphic.
\end{definition}
Holomorphic functions are automatically smooth. Also note that the transition functions are holomorphic \emph{if and only if} $\mathrm{d} (\text{transition} )\in \operatorname{GL} (n,\mathbb{C} )$.

Consider $\operatorname{SL} (n,\mathbb{R} )\subseteq \operatorname{GL} (n,\mathbb{R} )$ of matrices with determinant 1. In particular,
\[
	A^*\left( \mathrm{d} x^1 \wedge \cdots \wedge \mathrm{d} x^n \right) =\det A\left( \mathrm{d} x^1 \wedge \cdots \wedge \mathrm{d} x^n \right) =\mathrm{d} x^1 \wedge \cdots \wedge \mathrm{d} x^n
.\]
Thus $A$ preserves the volume form.

We can also consider $\mathrm{O} (n,\mathbb{R} )\subseteq \operatorname{GL} (n,\mathbb{R} )$ defined by $A^tA=1$. We also have $\operatorname{SO}(n) =\operatorname{SL}(n,\mathbb{R} ) \cap \operatorname{O}(n,\mathbb{R} ) $. We also have $\operatorname{U} (n,\mathbb{C} )$ defined as the set of matrices with $UU^\dagger=1$, and $\operatorname{SU} (n)=\operatorname{SL} (n,\mathbb{C} )\cap \operatorname{U} (n,\mathbb{C} )$.

It is at this point relevant to recall the Riemannian metric $g(x,y)=x\cdot y=x^ty$ on $\mathbb{R} ^n$ and the Riemannian metric $g(x,y)=x^\dagger y$ on $\mathbb{C} ^n$. We usually write, for $v,w\in\mathbb{C} ^n$, with $v=a+bi$ and $w=c+di$,
\[
	v^\dagger w = g\left( \begin{pmatrix}
	a \\
	b
	\end{pmatrix},\begin{pmatrix}
	c \\
	d
	\end{pmatrix} \right) +\omega \left( \begin{pmatrix}
	a \\
	b
	\end{pmatrix},\begin{pmatrix}
	c \\
	d
	\end{pmatrix}\right) 
\]
where $\omega (v,w)$ is defined by
\[
	\omega (v,w)=\begin{pmatrix}
		a^t & b^t
	\end{pmatrix}\begin{pmatrix}
	0 & -I \\
	-I & 0
	\end{pmatrix}\begin{pmatrix}
	c \\
	d
	\end{pmatrix}
.\]
We call $\omega $ a \emph{symplectic form}.

These structures are related by $\omega (v,w)=g(Jv,w)$, and thus if a map preerves the metric and the (complex) manifold structure, it preserves the Hermitian inner product and symplectic structure.

\begin{definition}[\cite{huybrechts}]\label{1.1.5}
	Let $V$ be a vector space. A \emph{almost complex structure} on $V$ is a map $J : V \to V$ with $J^2=-1$. An almost complex structure on an inner product space is \emph{compatible} with the inner product if $\left\langle J(-),J(-) \right\rangle =\left\langle -,- \right\rangle $. An \emph{almost complex structure} on a manifold $X$ is a vector field $J\in \mathfrak{X} (X)$ such that for all $x\in X$, the map $J_x$ is an almost complex structure on $T_xX$.
\end{definition}


\begin{definition}[\cite{huybrechts}]\label{1.1.6}
	Let $X$ be a complex manifold. A Riemannian metric $g$ on $X$ is a \emph{hermitian structure} on $X$ if for any $x\in X$, the inner product $g_x : T_xX\otimes T^*_xX \to \mathbb{R} $ is compatible with the complex structure $J_x$. The $(1,1)$-form $\omega \coloneqq g(J(-),-)$ is then called the \emph{fundamental form}.
\end{definition}

Let $g$ be the canonical Riemannian metric on $\mathbb{R} ^n$. On $\mathbb{C} ^n$, we can write the Riemannian metric $h(x,y)=x^\dagger y$ as follows. let $x=a+bi$ and $y=c+di$. We can interpret $x=(a,b)\in\mathbb{R} ^{2n} $ and similarly for $y$ in order to define $g(x,y)$. We can thus write
\[
	h(x,y)=g(x,y)+\omega (x,y),
\]
\[
	 \omega (x,y)=\begin{pmatrix}
	a^t & b^t
	\end{pmatrix}\begin{pmatrix}
	0 & -I \\
	-I & 0
	\end{pmatrix}\begin{pmatrix}
	c \\
	d
	\end{pmatrix}=\begin{pmatrix}
	a^t & b^t
	\end{pmatrix}\begin{pmatrix}
	-d \\
	-c
	\end{pmatrix}=-a^t d - b^tc
.\]
Note that $\omega $ is a symplectic form.

\begin{definition}[\cite{lee} 22]\label{1.1.6}
	Let $M$ be a manifold. A \emph{symplectic form} $\omega $ on $M$ is a 2-form such that for all $p\in M$, the vector $\omega _p$ is a nondegenerate 2-covector, meaning for all $v\in T_pM$ there exists $w\in T_pM$ such that $\omega_p(v,w)$.
\end{definition}

Note that $\omega (x,y)=g(J(x),y)$. Thus a map preserves $h$ and the complex manifold structure if and only if it preserves the Hermitian inner product and symplectic structure.
